\section{Introduction}

Machine translation for historical and ultra-low-resource languages is difficult for two separate reasons. First, parallel data is often scarce or nearly absent. Second, the outputs that matter to scholars are not always long, natural modern sentences: they may be titles, glosses, liturgical fragments, or compact bibliographic strings. Recent work on ancient-language NLP and historical-text translation has made this setting more visible, but it has also shown that data scarcity and evaluation instability are especially severe outside modern benchmark conditions \cite{sommerschield2023ancient, wang2023evahan, domingo2018modernizing}. Tangut is a canonical example. The script is historically important, but usable parallel Tangut-Chinese resources remain tiny. In the current repository snapshot, we only have 491 real Tangut-Chinese pairs.

This small-data regime also creates an obvious objection. Perhaps there is no need for task-specific adaptation at all: maybe a sufficiently strong Chinese frontier LLM, given dictionary glosses and a carefully engineered chain-of-thought-style prompt, can already do the task well enough. That objection is reasonable, so the paper should answer it directly rather than rhetorically. We therefore compare four families of methods in one frame: local dictionary-augmented prompting, a stronger proprietary frontier prompt-only baseline, synthetic pseudo-parallel SFT, and preference optimization on top of SFT \cite{chen2017incorporating, sennrich2016monolingual, rafailov2023dpo}. This turns the paper into a direct test of what prompting already buys, what local training still buys, and how evaluation changes the answer.

Our main finding is now three-part. First, the stronger prompt-only objection is real: the new DeepSeek-V3.2 few-shot dictionary baseline is the strongest non-reference system under the reference-aware judge, reaching 2.80/5 overall with 12 exact matches. Second, within the open/local trainable regime, structured synthetic supervision is still the strongest pure SFT recipe, and the strongest local systems appear only after preference optimization is moved onto that stronger base and the preference pairs are filtered much more aggressively. In particular, the strict gap-filtered DPO sigmoid variant reaches the highest local chrF++ in the repository (30.01), ties for the best local reference-aware overall score (2.22/5), keeps 5 exact matches, and stays close to reference title length. Third, the frontier comparison reveals complementary failure modes rather than simple dominance: the frontier model wins clearly on many canonical titles, but the local model still beats it on raw chrF++ and on a non-trivial minority of short, elliptical, or lexically awkward cases.

The repository also reveals two important failure modes. First, reference-free metrics can be badly misleading. The human reference itself is penalized by lexical coverage and perplexity, which means those metrics cannot serve as primary model-selection criteria for this task. Second, preference optimization is highly conditional. Legacy DPO on a weaker base underperforms badly, and even multitask-base DPO remains unstable when low-gap preference pairs are retained. The positive result therefore depends on both the base model and the pair-quality filter, not on DPO in the abstract.

The current draft therefore makes three contributions:
\begin{enumerate}
\item an empirical comparison of local prompting, frontier prompting, synthetic SFT, and DPO for Tangut short-text translation under only 491 real pairs;
\item a positive local-method result showing that structured multitask SFT is the right base and that strict gap-filtered DPO produces the strongest trainable local model, even exceeding the frontier prompt-only upper bound on chrF++;
\item a reference-aware evaluation and alignment analysis showing both why reference-free metrics and low-gap preference data can fail and why frontier prompting and local adaptation fail differently.
\end{enumerate}

The rest of the paper situates the task against related work, describes the experimental setting and the four comparison families, and then separates the positive local result from the frontier comparison and the metric/alignment analysis that makes the result interpretable.
